\section{Diagram 7}\label{sec:d7}

Diagram~7 is the butterfly of \cite[\S7.4]{AK} without circles.  Albouy and Kaloshin derive for
it the mass relations (36), (39) and (40) in \cite[\S\S8.7.1--8.7.3]{AK}.  We use none of these
relations.

\begin{theorem}\label{thm:d7}
Let $m\in\Rp^5$.  No Diagram~7 germ with an unbounded four-product lies on a curve through a
real point: if an irreducible curve $\Gamma\subset A_m$ carries such a germ, then $\Gamma$
contains no real point.
\end{theorem}

This is part~(b) of Theorem~\ref{thm:exclusion} for $r=7$.  Germs with $a+b\ge0$ in
Proposition~\ref{d7-cover} have no unbounded four-product (Lemma~\ref{d7-bounded}); the others
lie in the open cone of a nested chart, where an identity valid modulo \eqref{eq:system} gives
$I\to0$ (Lemma~\ref{d7-vanish}), so $I=U\equiv0$ on $\Gamma$, which is impossible at a real point.

\subsection{The diagram and the nested chart}
In the numbering of \cite[\S8.7]{AK}, the \emph{upper wing} $\{1,2,5\}$ is an isolated triangle
of $z$-edges without $z$-circles, the \emph{lower wing} $\{3,4,5\}$ is an isolated triangle of
$w$-edges without $w$-circles, and the four \emph{crossed} pairs $kj$ ($k\in\{1,2\}$,
$j\in\{3,4\}$) are not edges.
Since Diagram~7 has no circles, Rule~1e shows that neither wing has a maximal stroke (see the
proof of Proposition~\ref{d7-cover}(i)).  By Definition~\ref{def:germ} and the convention after
Theorem~\ref{thm:exclusion}, we may use this numbering and colouring.

\emph{Orders.}  As in Lemma~\ref{lem:germ} we renormalize a germ so that $Z$ and $W$ have a
common order, and reparametrize it by a power $t$ of its Puiseux parameter so that this order is
$-2$; thus $t$ plays the role of $\epsilon$ in \cite[\S3]{AK}, and $\ord$ is the order in $t$.
Renormalizations $(z,w)\mapsto(\lambda z,\lambda^{-1}w)$ preserve the diagram, the ties, $I$, $U$
and every $r_{kl}^2=z_{kl}w_{kl}$.

\emph{The nested chart} with exponents $P,Q,R$ has the coordinates $A_1,A_2,C,u_3,u_4$,
$B_3,B_4,D,v_1,v_2$ and $D_{kl}$ ($k<l$):
\begin{gather*}
  z_1=A_1,\quad z_2=A_2,\qquad w_3=t^{-Q}B_3,\quad w_4=t^{-Q}B_4,\\
  z_j=C+t^Pm_5u_j\ \ (j=3,4),\qquad z_5=C-t^P(m_3u_3+m_4u_4),\\
  w_k=t^{-Q}D+t^Rm_5v_k\ \ (k=1,2),\qquad w_5=t^{-Q}D-t^R(m_1v_1+m_2v_2),
\end{gather*}
and $\delta_{kl}=t^{e_{kl}}D_{kl}$, where $e_{kl}=-R/2$ on the upper wing, $Q/2$ on the crossed
pairs and $(Q-P)/2$ on the lower wing ($D$ without indices is not an inverse distance).
So $C$ and $t^{-Q}D$ are the centres of mass of the $z$-cluster $\{3,4,5\}$ and of the
$w$-cluster $\{1,2,5\}$, $P$ and $R$ are their internal orders, and $-Q$ is the order of the
common $w$-scale.
The exponents $e_{kl}$ balance all ten ties $\delta_{kl}^2z_{kl}w_{kl}=1$; those of the pairs
$12$ and $34$ read exactly
\[
  m_5D_{12}^2(A_2-A_1)(v_2-v_1)=1,\qquad m_5D_{34}^2(u_4-u_3)(B_4-B_3)=1.
\]
The \emph{open cone} is $P>Q>0$, $R>0$.

\emph{Gauges.}  The chart does not impose the gauge $z_{12}=w_{12}$.  By the displayed ties,
$A_2-A_1$ and $B_4-B_3$ never vanish, so they could serve as gauges for the rotations and for
the rescalings of $t$.  We fix none, since $I$, $U$, the ties and the zero set of the rows of
\eqref{eq:system} are invariant under the rotations.

\subsection{Coverage}

\begin{proposition}[Coverage]\label{d7-cover}
Let $\gamma$ be a Diagram~7 germ.  There are $a,b$ with $-2<a<1$ and $-2<b<1$ such that the
following hold.
\begin{romanlist}
  \item On the upper wing $\ord z_{kl}=a$ and $\ord w_{kl}=(4-a)/3$; on the lower wing
  $\ord w_{kl}=b$ and $\ord z_{kl}=(4-b)/3$; on the crossed pairs $\ord z_{kl}=a$ and
  $\ord w_{kl}=b$.  Moreover $\ord z_k\ge a$ and $\ord w_k\ge b$ for every $k$.
  \item After the rotation $(z,w)\mapsto(t^{-a}z,t^{a}w)$, $\gamma$ is a germ of the nested
  chart with exponents
  \[
    P=\frac{4-b}{3}-a,\qquad Q=-(a+b),\qquad R=\frac{4+2a}{3};
  \]
  all chart coordinates have nonnegative order, so they converge, and every $D_{kl}$ has order
  zero.  Moreover $R=\frac23(a+2)>0$ and $P-Q=\frac23(b+2)>0$, so $(P,Q,R)$ lies in the open
  cone if and only if $a+b<0$.
\end{romanlist}
\end{proposition}

\begin{proof}
(i) We use Fact~\ref{fact:rules}.  By Rule~2e \cite[p.~552]{AK} the three $z$-separations of a
wing have one order, and so do its three $w$-separations; let $a$ be the $z$-order of the upper
wing and $b$ the $w$-order of the lower wing.  As $\delta_{kl}^2z_{kl}w_{kl}=1$,
$\ord Z_{kl}=-(\ord z_{kl}+3\ord w_{kl})/2$ and $\ord W_{kl}=-(3\ord z_{kl}+\ord w_{kl})/2$;
since the strokes have order $-2$, this gives the other two wing orders.  The upper wing is an
isolated component of the $z$-diagram without $z$-circles, so it has no maximal $z$-stroke
(Rule~1e, p.~546), and along a $z$-stroke Estimate~2 (p.~548) gives $-2\le\ord z_{kl}\le1$, with
$-2$ exactly for a maximal $z$-edge and $1$ exactly for a $zw$-edge.  Hence $-2<a<1$, and dually
$-2<b<1$.  As $z_{13}=z_{15}+z_{53}$
with $\ord z_{53}=(4-b)/3>1>a$, and $w_{13}=w_{15}+w_{53}$ with $\ord w_{15}=(4-a)/3>1>b$, the
crossed orders are $a$ and $b$; likewise for $14$, $23$, $24$.  Finally
$\sum_km_kz_k=\sum_km_kw_k=0$ on $A_m$ (proof of Lemma~\ref{lem:distances}), so each $z_k$ is a
combination of $z$-differences, all of order at least $a$; likewise $\ord w_k\ge b$.

(ii) After the rotation the $z$-differences have order $0$ on the upper wing and the crossed
pairs and $P$ on the lower wing, the $w$-differences have order $R$ on the upper wing and $-Q$
elsewhere, and $\ord z_k\ge0$, $\ord w_k\ge-Q$.  Inverting the chart map, $A_k=z_k$,
$B_j=t^Qw_j$, $C$ and $t^{-Q}D$ are the centres of the two clusters, $m_5t^Pu_j=z_j-C$ and
$m_5t^Rv_k=w_k-t^{-Q}D$ are combinations of differences inside them, and
$D_{kl}=t^{-e_{kl}}\delta_{kl}$, where $\ord\delta_{kl}=-(\ord z_{kl}+\ord w_{kl})/2=e_{kl}$.  So
every coordinate has nonnegative order and $\ord D_{kl}=0$.
\end{proof}

Thus a germ lies in the open cone exactly when $a+b<0$, and the germs with $a+b\ge0$ have no
unbounded four-product (Lemma~\ref{d7-bounded}).  In the terminology of \cite[\S8.7]{AK}, a germ
with a wing in the general case has $a+b<0$; we shall not need this.

\subsection{The inertia identity}
Write
\[
  F^z_k=z_k-\sum_{l\ne k}m_l\delta_{kl}^3(z_k-z_l),\qquad
  F^w_k=w_k-\sum_{l\ne k}m_l\delta_{kl}^3(w_k-w_l)
\]
for the coordinate rows of \eqref{eq:system} and $T_{kl}=\delta_{kl}^2z_{kl}w_{kl}-1$ for the
ties.  Then $I-U=\sum_km_kw_kF^z_k+\sum_{k<l}m_km_l\delta_{kl}T_{kl}$, so $U=I$ on $A_m$
\cite[p.~541]{AK}.

\begin{lemma}[Inertia identity]\label{d7-identity}
Let
\[
  S_5=\sum_{j=3,4}\sum_{l=1,2}m_jm_l\delta_{jl}^3
  \bigl((z_j-z_5)(w_l-w_5)-(w_j-w_5)(z_l-z_5)\bigr).
\]
Then, as polynomials in the masses and in $(z,w,\delta)$,
\begin{multline*}
  I=\sum_{k=1,2}m_kz_k(w_k-w_5)+\sum_{j=3,4}m_jw_j(z_j-z_5)+S_5\\
  -\sum_{j=3,4}m_j\bigl((z_j-z_5)F^w_j-(w_j-w_5)F^z_j\bigr)+w_5\sum_{k=1}^5m_kF^z_k.
\end{multline*}
In particular $I=\sum_{k=1,2}m_kz_k(w_k-w_5)+\sum_{j=3,4}m_jw_j(z_j-z_5)+S_5$ on $A_m$.
\end{lemma}

\begin{proof}
Put $\hat z_k=z_k-z_5$ and $\hat w_k=w_k-w_5$.  By antisymmetry of the force terms,
$\sum_km_kF^z_k=\sum_km_kz_k$.  Since $z_j\hat w_j-w_j\hat z_j=z_5w_j-w_5z_j$,
\begin{align*}
  I&=\sum_km_kz_k\hat w_k+w_5\sum_km_kz_k\\
   &=\sum_{k=1,2}m_kz_k\hat w_k+\sum_{j=3,4}m_j\bigl(w_j\hat z_j+z_5w_j-w_5z_j\bigr)
    +w_5\sum_km_kF^z_k.
\end{align*}
For $j\in\{3,4\}$, using $z_j-z_l=\hat z_j-\hat z_l$ and $w_j-w_l=\hat w_j-\hat w_l$,
\[
  \hat z_jF^w_j-\hat w_jF^z_j=(w_5z_j-z_5w_j)
  +\sum_{l\ne j}m_l\delta_{jl}^3(\hat z_j\hat w_l-\hat w_j\hat z_l).
\]
In $\sum_{j=3,4}m_j(\hat z_jF^w_j-\hat w_jF^z_j)$ the terms with $l=5$ vanish, those with
$l\in\{3,4\}$ cancel in pairs, and those with $l\in\{1,2\}$ add up to $S_5$.  Hence
$\sum_{j=3,4}m_j(z_5w_j-w_5z_j)=S_5-\sum_{j=3,4}m_j(\hat z_jF^w_j-\hat w_jF^z_j)$, and
substituting this into the display gives the identity.
\end{proof}

The term $S_5$ is the torque about body~$5$ of the crossed forces on the lower wing; compare
\cite[(31)]{AK}.

\begin{lemma}\label{d7-vanish}
Let $\gamma$ be a Diagram~7 germ on $A_m$ with $a+b<0$ (Proposition~\ref{d7-cover}).  Then
$I(\gamma(t))\to0$ as $t\to0$.
More precisely, in the nested chart of Proposition~\ref{d7-cover}(ii) the three terms
$\sum_{k=1,2}m_kz_k(w_k-w_5)$, $\sum_{j=3,4}m_jw_j(z_j-z_5)$ and $S_5$ are sums of powers $t^e$
times polynomials in the chart coordinates, with exponents
\[
  e=R;\qquad e=P-Q;\qquad e\in\bigl\{\tfrac Q2,\ P+\tfrac Q2,\ \tfrac{3Q}2+R,\ P+\tfrac{3Q}2+R\bigr\},
\]
respectively.
\end{lemma}

\begin{proof}
On $A_m$ the last two sums in Lemma~\ref{d7-identity} vanish, and the three remaining terms are
invariant under the rotations, so we may evaluate them on the rotated germ in the chart.  There
\begin{gather*}
  w_k-w_5=t^R\bigl((m_k+m_5)v_k+m_{k'}v_{k'}\bigr),\qquad
  z_j-z_5=t^P\bigl((m_j+m_5)u_j+m_{j'}u_{j'}\bigr),\\
  w_j-w_5=t^{-Q}(B_j-D)+t^R(m_1v_1+m_2v_2),\\
  z_l-z_5=A_l-C+t^P(m_3u_3+m_4u_4),
\end{gather*}
where $\{k,k'\}=\{1,2\}$, $\{j,j'\}=\{3,4\}$ and $l\in\{1,2\}$, and
$\delta_{jl}^3=t^{3Q/2}D_{jl}^3$.  Substituting, together with $z_k=A_k$ and $w_j=t^{-Q}B_j$,
gives the stated exponents.  Since $a+b<0$, Proposition~\ref{d7-cover}(ii) puts $(P,Q,R)$ in the
open cone, so every exponent is positive, and all chart coordinates are bounded.  Hence the three
terms, and with them $I$, tend to zero.
\end{proof}

\subsection{Germs outside the open cone}

\begin{lemma}\label{d7-bounded}
Let $\gamma$ be a Diagram~7 germ with $a+b\ge0$ (Proposition~\ref{d7-cover}).  Then every
four-product tends to zero along $\gamma$.  If $a+b=0$, then $-1<a<1$, and $p_{12}$, $p_{34}$
and each of the other eight four-products have the orders $(16-4a)/3$, $(16+4a)/3$ and $8/3$.
If $a+b>0$, every $r_{kl}$ tends to zero, and so does $I$.
\end{lemma}

\begin{proof}
By Proposition~\ref{d7-cover}(i), $r_{kl}^2=z_{kl}w_{kl}$ has order $(4+2a)/3>0$ on the upper
wing, $(4+2b)/3>0$ on the lower wing and $a+b\ge0$ on the crossed pairs.  The four-product
$p_{ij}$ contains the three squared sides of the triangle $\{k,l,m\}$ complementary to
$\{i,j\}$, and this triangle contains a wing pair: if $5\in\{k,l,m\}$, both pairs through~$5$
are wing pairs; otherwise $\{k,l,m\}\subset\{1,2,3,4\}$ contains $\{1,2\}$ or $\{3,4\}$ (the
crossed pairs form the bipartite graph $K_{2,2}$).  Hence
$\ord p_{ij}\ge\min\{(4+2a)/3,(4+2b)/3\}>0$.  If $a+b=0$, the bounds $-2<a,b<1$ give $-1<a<1$,
and adding orders gives the stated values.  If $a+b>0$, every $r_{kl}^2$ has positive order, and
the Lagrange identity $MI=(\sum_km_kz_k)(\sum_km_kw_k)+\sum_{k<l}m_km_lz_{kl}w_{kl}$, with
$\sum_km_kz_k=\sum_km_kw_k=0$, gives $I\to0$.
\end{proof}

\begin{remark}\label{d7-balanced}
The germs with $a+b=0$ are not excluded, and need not be: they carry no unbounded four-product,
so Proposition~\ref{prop:reduction} never produces them.
\end{remark}

\subsection{End of the proof}

\begin{proof}[Proof of Theorem~\ref{thm:d7}]
Let $\gamma$ be a Diagram~7 germ with an unbounded four-product on an irreducible curve
$\Gamma\subset A_m$, numbered as above, and let $a,b$ be as in Proposition~\ref{d7-cover}.  Some
four-product has negative order along $\gamma$ (Definition~\ref{def:germ}), so $a+b<0$ by
Lemma~\ref{d7-bounded}, and $I(\gamma(t))\to0$ by Lemma~\ref{d7-vanish}.  By
Lemma~\ref{lem:U-constant}, $I=U$ is constant on $\Gamma$, with value
$\lim_{t\to0}I(\gamma(t))=0$.  At a real point $w_k=\bar z_k$, so $I=\sum_km_k|z_k|^2$,
which is positive because $m_k>0$ and the tie of the pair $12$ forces $z_1\ne z_2$.  Hence
$\Gamma$ contains no real point.
\end{proof}

\begin{remark}\label{d7-masses}
The proof of Theorem~\ref{thm:d7} uses no relation among the masses: Lemma~\ref{d7-identity} is a
polynomial identity, Proposition~\ref{d7-cover} and Lemma~\ref{d7-bounded} use only Rules~1e,
2e and Estimate~2 of \cite[\S3]{AK}, Lemma~\ref{d7-vanish} only the signs of the exponents, and
the last step only Lemma~\ref{lem:U-constant} and $m_k>0$.  So the theorem does not depend on
whether the masses satisfy (36), (39) or (40), and no mass condition is derived for Diagram~7.
\end{remark}
